\chapter{Part 2a --- 32-bit Hierarchical ALU}
\label{ch:part2a}

The first 32-bit realisation is \emph{structural}: it instantiates
thirty-two copies of the verified 1-bit slice and wires their carries into
a ripple chain. This is the direct expression of the bit-slice methodology
and is the design ultimately chosen for the chip.

\section{Bit-slice chaining}
\Cref{fig:chain} shows the construction. Slice~$i$ takes operand bits
$A_i$, $B_i$ and the carry $C_i$ from the slice below, and produces result
bit $F_i$ and the carry $C_{i+1}$ for the slice above. The external
\bus{Cin} enters slice~0; the final carry $C_{32}$ becomes \bus{Cout}. The
same \sel{} bus fans out to every slice, so all slices perform the same
operation on their respective bit.

\begin{figure}[htbp]
\centering
\begin{tikzpicture}[font=\small,node distance=6mm]
  \foreach \i/\x in {0/0, 1/2.6, 2/5.2}{
    \node[arith,minimum width=18mm,minimum height=14mm] (s\i) at (\x,0)
         {ALU\\ slice \i};
  }
  \node at (7.0,0) {$\cdots$};
  \node[arith,minimum width=18mm,minimum height=14mm] (s31) at (8.8,0)
       {ALU\\ slice 31};

  % Carry chain (left to right)
  \draw[sig] (-1.6,0) node[left,buslbl]{Cin} -- (s0.west);
  \draw[sig] (s0.east) -- node[above,buslbl]{$C_1$} (s1.west);
  \draw[sig] (s1.east) -- node[above,buslbl]{$C_2$} (s2.west);
  \draw[sig] (s2.east) -- node[above,buslbl]{$C_3$} (6.6,0);
  \draw[sig] (7.4,0) -- node[above,buslbl]{$C_{31}$} (s31.west);
  \draw[sig] (s31.east) -- (10.4,0) node[right,buslbl]{Cout};

  % Operand inputs (top) and results (bottom)
  \foreach \i in {0,1,2,31}{
    \draw[sig] ($(s\i.north)+(-0.35,0.7)$) node[above,buslbl]{$A_{\i}$}
          -- ($(s\i.north)+(-0.35,0)$);
    \draw[sig] ($(s\i.north)+(0.35,0.7)$) node[above,buslbl]{$B_{\i}$}
          -- ($(s\i.north)+(0.35,0)$);
    \draw[sig] (s\i.south) -- ($(s\i.south)+(0,-0.7)$)
          node[below,buslbl]{$F_{\i}$};
  }
\end{tikzpicture}
\caption{Ripple-carry chaining of thirty-two 1-bit ALU slices. The carry
propagates left-to-right from \bus{Cin} to \bus{Cout}; \sel{} (not shown)
broadcasts to all slices.}
\label{fig:chain}
\end{figure}

\lstinputlisting[style=vlog,firstline=1,lastline=14,
  caption={\texttt{alu\_32bit.v} --- port list, carry vector and the first
  slice instance (slices 1--31 are identical, indices incremented).},
  label={lst:alu32}]{rtl/alu_32bit.v}

The shift operations are handled once at the top level rather than inside
each slice: when \code{S3\,S2}$=10$ or $11$ the 32-bit result is
\code{a\_i>>1} or \code{a\_i<<1} and \bus{Cout} is forced to $0$, exactly
mirroring the per-slice output mux of \cref{fig:alu1top} but applied to the
whole word.

\paragraph{Design decision --- ripple carry.}
Chaining the slice carries produces a \emph{ripple-carry adder}: simple,
compact, and a perfect fit for the bit-slice structure, but with a carry
that must propagate through all 32 stages in the worst case. As the timing
results below confirm, this carry chain is the critical path of the
design --- the classic area-versus-speed trade-off of ripple carry against
carry-lookahead.

\section{Functional simulation}
The 32-bit structural ALU is simulated with wide operands across all
operation classes; \cref{fig:hiersim} shows the result matching the
expected arithmetic, logic and shift outputs.

\reportfig{0.95}{fig12_hier_sim}{32-bit hierarchical ALU simulation, built
from thirty-two 1-bit slices.}{fig:hiersim}

\section{Synthesis}
Genus synthesises the structural design to \textbf{367 cells}
(\cref{fig:hiersynth}). \Cref{tab:hierarea} summarises the area; the
thirty-two \texttt{ADDFX1} full adders ($\SI{164.2}{\micro\meter\squared}$)
are the arithmetic backbone, and clock/DRC buffers again make up a large
share.

\reportfig{0.95}{fig13_hier_synth}{Genus synthesis schematic of the 32-bit
hierarchical ALU (367 leaf cells, 32 slice modules).}{fig:hiersynth}

\begin{table}[htbp]
\centering
\caption{32-bit hierarchical ALU area breakdown (Genus).}
\label{tab:hierarea}
\small
\begin{tabular}{@{}l r r r@{}}
\toprule
Type & Instances & Area (\umsq) & Share \\
\midrule
Inverter & 35  & 24.282  & \SI{2.5}{\percent}\\
Buffer   & 65  & 325.584 & \SI{33.4}{\percent}\\
Logic    & 267 & 625.586 & \SI{64.1}{\percent}\\
\midrule
\textbf{Total} & \textbf{367} & \textbf{975.452} & \SI{100}{\percent}\\
\bottomrule
\end{tabular}
\end{table}

\reportfig{0.55}{fig14_hier_area}{Genus area report, 32-bit hierarchical
ALU.}{fig:hierarea}

\section{Power}
Estimated total power is \textbf{\SI{25.32}{\micro\watt}}
(\cref{fig:hierpower}), again dominated by switching activity.

\reportfig{0.75}{fig15_hier_power}{Genus power report, 32-bit hierarchical
ALU.}{fig:hierpower}

\section{Place and route}
\Cref{fig:hierpnr} shows the placed-and-routed core in Innovus.

\reportfig{0.6}{fig16_hier_pnr}{Innovus place-and-route of the 32-bit
hierarchical ALU.}{fig:hierpnr}

\section{Physical verification}
\begin{figure}[htbp]
\centering
\begin{subfigure}{0.48\linewidth}\centering
  \includegraphics[width=\linewidth]{fig17_hier_drc}
  \caption{DRC results}\label{fig:hierdrc}
\end{subfigure}\hfill
\begin{subfigure}{0.48\linewidth}\centering
  \includegraphics[width=\linewidth]{fig18_hier_conn}
  \caption{Connectivity check}\label{fig:hierconn}
\end{subfigure}
\caption{Physical signoff of the 32-bit hierarchical ALU: DRC and
connectivity both clean.}
\label{fig:hiersign}
\end{figure}

\section{Timing and maximum speed}
Static timing analysis reports a worst-case critical path of
\textbf{\SI{9.724}{\nano\second}}. As anticipated, the path runs through
the \textbf{ripple-carry chain}: from an early select/operand bit, through
the arithmetic mux of slice~0, and then \code{CI}$\rightarrow$\code{CO}
across all thirty-two \texttt{ADDFX1} adders to \code{f\_o[31]}. Each adder
stage contributes roughly \SI{186}{\pico\second} of carry delay, so the
32-stage chain dominates. The estimated maximum frequency is
\begin{equation}
  f_{\max} \approx \frac{1}{\SI{9.724}{\nano\second}}
           \approx \SI{102.8}{\mega\hertz}.
\end{equation}
This ripple-carry delay is the price of the simple, highly regular
bit-slice structure, and it is the number carried into the comparison of
\cref{ch:compare}.
